\begin{resumo}
O Sistema Integrado de Gestão de Atividades Acadêmicas (SIGAA) é o ponto único de acesso do estudante da Universidade de Brasília (UnB) a notas, horários, histórico, documentos e turmas virtuais, sem alternativa institucional. Relatos recorrentes de dificuldade de uso, a interface pensada para computador de mesa e a existência de ferramentas comunitárias que suprem funções do sistema motivam uma avaliação objetiva da sua qualidade. Este artigo apresenta a especificação dessa avaliação, conduzida pela equipe Adele Goldstine segundo o processo de avaliação da família SQuaRE (ISO/IEC 25040) e o modelo de qualidade de produto da ISO/IEC 25010:2023. Na Fase 1, foram definidos o propósito, as partes interessadas, o modelo de qualidade adaptado e o escopo, e uma matriz Impacto~$\times$~Risco com filtro de viabilidade selecionou quatro características: Capacidade de Interação (Usabilidade), Adequação Funcional, Eficiência de Desempenho e Flexibilidade (Portabilidade). Na Fase 2, o método \textit{Goal-Question-Metric} (GQM) desdobrou essas características em quatro objetivos, 14 questões com hipóteses e 20 métricas, cada uma com procedimento de coleta, níveis de pontuação e critérios de julgamento definidos antes da coleta. A especificação combina teste com usuários, inspeção por dois avaliadores, ferramentas automáticas e um questionário que compara estudantes de diferentes semestres. A gestão do projeto é acompanhada por indicadores do PSM/CID. As fases seguintes executarão a medição e uma ação de melhoria baseada nos resultados.

\vspace{\onelineskip}
\noindent
\textbf{Palavras-chave}: \imprimirpalavrachaveum. \imprimirpalavrachavedois. \imprimirpalavrachavetres. \imprimirpalavrachavequatro.
\end{resumo}
